\documentclass[12pt,a4paper]{article}

\usepackage[margin=1in]{geometry}
\usepackage{amsmath, amssymb}
\usepackage{booktabs}
\usepackage{enumitem}
\usepackage{hyperref}
\usepackage{parskip}
\usepackage{titlesec}
\usepackage{array}
\usepackage{xcolor}
\usepackage{fancyhdr}
\setlength{\headheight}{14.5pt}
\addtolength{\topmargin}{-2.5pt}
\usepackage{graphicx}

\hypersetup{
    colorlinks=true,
    linkcolor=blue!60!black,
    urlcolor=blue!60!black,
}

\pagestyle{fancy}
\fancyhf{}
\rhead{Engineering Standards and Constraints}
\lhead{QUBO-Based DFL Node Placement}
\cfoot{\thepage}

\titleformat{\section}{\large\bfseries}{}{0em}{}[\titlerule]
\titleformat{\subsection}{\normalsize\bfseries}{}{0em}{}

\title{
    \textbf{Engineering Standards and Constraints} \\[0.5em]
    \large QUBO-Based Sensor Node Placement for Device-Free Localization
}
\author{Mohamed Khalil Brik}
\date{Spring 2026}

\begin{document}

\maketitle
\tableofcontents
\newpage

% ─────────────────────────────────────────────────────────────────────────────
\section{Introduction}

This document identifies the engineering standards and constraints that
govern the design, implementation, and evaluation of the QUBO-based sensor
node placement system. It serves as a formal record of the standards followed
and the boundaries within which the system was built, as required for the
final thesis submission package.

% ─────────────────────────────────────────────────────────────────────────────
\section{Applicable Engineering Standards}

\subsection{IEEE 830: Software Requirements Specification}

The System Requirements and Design Specifications document follows the
structure recommended by IEEE Std 830-1998. Requirements are written using
consistent identifiers (FR-1 through FR-10, NFR-1 through NFR-5), are
individually testable, and are organized by functional area. This ensures
the requirements are complete, unambiguous, and traceable to the
implementation.

\subsection{IEEE 829: Software Test Documentation}

The Verification and Validation Report follows the test documentation
practices of IEEE Std 829-2008. Each test case includes a unique identifier,
the requirement being verified, the input conditions, the expected outcome,
and the actual outcome. This structure makes the testing process auditable
and reproducible.

\subsection{IEEE 1012: Software Verification and Validation}

The overall V\&V approach is aligned with IEEE Std 1012-2016. Verification
checks that each component was built correctly according to the specification.
Validation checks that the full system produces results that are meaningful
and consistent with the research goals. Both levels are addressed in the
accompanying V\&V report.

\subsection{IEEE 754: Floating-Point Arithmetic}

All floating-point computations in the system rely on IEEE 754 double
precision (64-bit) arithmetic, which is the default in Python's
\texttt{numpy} and \texttt{pandas} libraries. The binary search tolerance
$\epsilon = 0.01$ is chosen to be well above the resolution of double
precision arithmetic, avoiding any numerical instability near the convergence
boundary.

\subsection{IEEE Publication Figure Standards}

All figures generated for the thesis follow the IEEE journal figure
guidelines. Figures use a serif font (Computer Modern), a minimum font size
of 9~pt, a line width of at least 0.8~pt, and are saved at 300~DPI. Colors
are chosen to remain distinguishable in both color and grayscale print.
Figure width is set to 3.5 inches to match a single IEEE journal column.

\subsection{Reproducible Research Standards}

The system follows the reproducibility principles outlined by ACM and IEEE
for computational research. All stochastic components use a fixed random
seed (seed = 42 for single-run experiments; seeds 0 through 99 for the
multi-run experiment). All input data and output results are stored as
plain CSV files so that any reviewer can re-run the pipeline and obtain
identical results.

% ─────────────────────────────────────────────────────────────────────────────
\section{Design Constraints}

\subsection{Hardware Constraints}

\begin{itemize}[leftmargin=2em]
    \item \textbf{No quantum hardware access.} The system was designed to run
          entirely on classical hardware. The quantum annealing solver is
          emulated using OpenJij's \texttt{SQASampler}, which simulates the
          behavior of a D-Wave processor. The projected hardware anneal time
          of 20~$\mu$s per read is used for TTS calculations, but no actual
          D-Wave machine was accessed.
    \item \textbf{No specialized GPU or FPGA.} The system runs on a standard
          laptop CPU. All matrix operations use NumPy, which may use
          multi-threaded BLAS internally but does not require a GPU.
    \item \textbf{Sensor node count fixed at $N = 10$.} The environment
          provides exactly 10 candidate nodes, yielding a finite set of
          AP-MP pairs. The system is not required to scale beyond this
          dataset for the thesis.
\end{itemize}

\subsection{Software Constraints}

\begin{itemize}[leftmargin=2em]
    \item \textbf{Python 3.8 or later.} The system uses f-strings,
          \texttt{pathlib}, and type hints that require at least Python 3.8.
    \item \textbf{OpenJij dependency.} The QUBO solvers depend on the OpenJij
          library. The system does not implement its own annealing algorithm.
          Any future change to the OpenJij API could require updates to the
          solver scripts.
    \item \textbf{No internet connection required at runtime.} All data is
          stored locally. The system does not call any external APIs or
          download any files during execution.
\end{itemize}

\subsection{Budget Constraints}

\begin{itemize}[leftmargin=2em]
    \item \textbf{AP budget: $n \leq 5$.} At most 5 distinct Access Point
          nodes may appear in the selected set of links.
    \item \textbf{MP budget: $m \leq 5$.} At most 5 distinct Monitoring
          Point nodes may appear in the selected set of links.
    \item These budgets reflect realistic hardware deployment limits and
          are enforced as hard constraints throughout the binary search.
          Any solution that violates either budget is immediately rejected.
\end{itemize}

\subsection{Algorithmic Constraints}

\begin{itemize}[leftmargin=2em]
    \item \textbf{Binary search tolerance $\epsilon = 0.01$.} The search
          terminates when the interval $[a, b]$ shrinks to less than 0.01.
          This results in at most $\lceil \log_2(1/0.01) \rceil = 7$
          iterations.
    \item \textbf{QUBO formulation is unconstrained.} Budget feasibility
          is not encoded as a penalty term inside the QUBO. This is a
          deliberate design choice to avoid penalty weight tuning. The QUBO
          only balances importance against redundancy through the single
          parameter $\alpha$.
    \item \textbf{k-NN with $k = 3$ fixed.} The localization evaluator uses
          exactly 3 nearest neighbors. This value was chosen to balance
          between noise sensitivity and spatial resolution and is held
          constant across all experiments for fair comparison.
\end{itemize}

\subsection{Data Constraints}

\begin{itemize}[leftmargin=2em]
    \item \textbf{Fixed dataset.} The fingerprint RSS data, coordinates,
          and precomputed importance and redundancy matrices are fixed inputs.
          The system does not perform any data collection or augmentation.
    \item \textbf{Train/test split is predetermined.} The split between
          training and test fingerprints is defined in the input CSV files.
          The system does not shuffle or re-split the data.
    \item \textbf{Importance scores are RSS variance.} The importance of
          each AP-MP link is defined as the variance of its RSS measurements
          across fingerprint locations. Links with higher variance carry more
          discriminative power for localization.
    \item \textbf{Redundancy is measured by RSS correlation.} The redundancy
          between two links is the Pearson correlation of their RSS vectors.
          Selecting two highly correlated links adds little new information.
\end{itemize}

\subsection{Time Constraints}

\begin{itemize}[leftmargin=2em]
    \item \textbf{Single-run binary search must complete within 10 minutes}
          on a standard laptop. In practice, SA with 100 reads and 1000
          sweeps per iteration converges in under 2 minutes.
    \item \textbf{The 100-trial multi-run experiment} is intended as an
          offline batch job and has no strict time limit, though it typically
          completes within 3 hours on the same hardware.
\end{itemize}

\subsection{Ethical and Academic Constraints}

\begin{itemize}[leftmargin=2em]
    \item \textbf{No human subject data.} The RSS fingerprint dataset does
          not contain any personally identifiable information. It consists
          only of signal strength measurements at fixed physical locations.
    \item \textbf{Reproducibility.} All experiments use fixed seeds so that
          claims made in the thesis can be independently verified from the
          submitted code and data.
    \item \textbf{Honest reporting.} The greedy baseline is included and
          honestly compared against the QUBO methods. No results are
          selectively omitted.
\end{itemize}

% ─────────────────────────────────────────────────────────────────────────────
\section{Summary Table}

\begin{table}[h!]
\centering
\renewcommand{\arraystretch}{1.3}
\begin{tabular}{@{}lp{4.5cm}p{6cm}@{}}
\toprule
\textbf{Standard / Constraint} & \textbf{Category} & \textbf{Impact on System} \\
\midrule
IEEE 830   & Documentation    & Structure of SRS/SDS document \\
IEEE 829   & Testing          & Structure of V\&V test cases \\
IEEE 1012  & V\&V process     & Scope of verification and validation \\
IEEE 754   & Arithmetic       & Floating-point stability of binary search \\
IEEE figures & Visualization  & DPI, font size, line width of all figures \\
Reproducibility & Research    & Fixed seeds, plain CSV outputs \\
AP budget $\leq 5$ & Hardware & Hard constraint in binary search \\
MP budget $\leq 5$ & Hardware & Hard constraint in binary search \\
$\epsilon = 0.01$ & Algorithmic & Binary search convergence bound \\
$k = 3$ (k-NN) & Algorithmic  & Fixed evaluation metric \\
No real QA hardware & Hardware & SQA emulation via OpenJij \\
Python 3.8+ & Software        & Runtime environment \\
No internet & Software        & Fully offline execution \\
No human data & Ethical       & No IRB approval required \\
\bottomrule
\end{tabular}
\caption{Summary of applicable standards and constraints.}
\end{table}

\end{document}
